% QUESTAO
Observe uma circunferência
\vspace{-1mm}
\begin{center}
\tikz[>=latex,scale=0.2,rotate=0]{
\pgfmathsetmacro{\R}{sqrt(52)};
\coordinate (O) at (0,0);
\fill (O) circle (0.4*\R mm);
\draw[thick] (O) circle (\R);	
\draw[thick] (-135:\R) -- (45:\R);
\draw[dashed] (-135:\R) -- ($(-135:\R)+(-45:\R+1)$) (45:\R) -- ($(45:\R)+(-45:\R+1)$);
\draw[ decorate, decoration={brace,raise=0pt} ] ($(45:\R)+(-45:\R+1)$) -- ($(-135:\R)+(-45:\R+1)$);
\node at (-45:10) {$D$};
\node[below] at (45:\R/1.7) {$R$};
\node[below] at (-135:\R/2.7) {$R$};
}
\end{center}
\vspace{-2mm}
Qual é a razão entre o raio $R$ e o diâmetro $D$?

% ALTERNATIVAS 3
$\frac{R}{D}=\frac{1}{2}$
$\frac{R}{D}=\frac{3}{2}$
$\frac{R}{D}=2$
$\frac{R}{D}=1$
$\frac{R}{D}=\frac{2}{3}$



